\section{Introduction}
\label{sec:introduction}

% Part 1: the core problem
Autism spectrum disorder (ASD) is recognised as a neurodevelopmental condition in which social communication and behavioural patterns are affected \citep{Maenneretal2023}. ASD diagnosis is dependent on behavioural
assessment by trained clinicians, supported by the Autism Diagnostic Observation Schedule (ADOS-2)
\citep{CHOP_ADOS2} and the Autism Diagnostic Interview--Revised \citep{Lordetal1994}.
Each assessment is conducted in person by trained clinicians through behavioural observation and clinical judgement, making the diagnostic process both time- and resource-intensive. Timely access to assessment is therefore important, as earlier ASD identification enables appropriate interventions and support during critical stages of development, in line with current clinical recommendations \citep{Zwaigenbaumetal2015}. Brain-based measurements are therefore being investigated as research tools for characterising autism-related variation and for studying whether diagnostic groups can be distinguished from neuroimaging data.


Resting-state functional MRI (rs-fMRI) provides estimates of functional connectivity from spontaneous blood-oxygen-level-dependent activity across brain regions. These connectomes have been widely used for autism-related group comparisons and classification. These signals are then used to estimate functional connectivity, which is defined as the statistical relationship between the activity of different brain regions. In autistic individuals, differences in functional connectivity are observed across several brain regions and networks in large multisite studies. Functional connectivity therefore serves as an input for group comparisons and ASD classification. Large multisite resources such as the Autism Brain Imaging Data Exchange I (ABIDE-I) have enabled connectome-based autism classification to be studied across heterogeneous acquisition sites and participant populations. In ABIDE-I, rs-fMRI and phenotypic data from 1,112 participants across multiple sites are made openly available \citep{DiMartinoetal2014}. The dataset is based on a research population, in which participants are recruited for research studies and most individuals in the non-autistic group are typically developing controls.


% Part 2: the problem identified
A central problem in connectome-based autism classification is that measured performance and biological interpretation can be entangled with cohort composition, motion, feature construction, regularisation, and model design. ABIDE-I provides a suitable multisite setting in which these factors can be examined within the same analytical framework. Systematic reviews of ABIDE-I studies report differences in sample selection, preprocessing, validation methods, and classifier design, which limit direct comparison across studies \citep{Santanaetal2022,Hudaetal2024}. Participant composition is one important source of variation. Cognitive ability differs substantially across autistic participants, yet classification performance is usually reported only for the overall cohort. In published ABIDE-I benchmarks, classification performance is generally reported for the overall cohort, and performance across different levels of cognitive ability is rarely examined \citep{Heinsfeldetal2018,Eslamietal2019,AlmuqhimSaeed2021}. Full-scale intelligence quotient (IQ) varies across autistic participants and is generally lower than that of non-autistic participants in the dataset (Table~\ref{tab:participants}). Classification performance may therefore be affected by the cognitive characteristics of the participants included in the analysis. Variation in classifier scores with verbal IQ has also been reported in ABIDE data \citep{Nielsenetal2013}. This evidence is not sufficient to establish whether the relationship is specifically associated with autism. Head motion creates a related interpretive problem because it alters measured functional connectivity and differs between diagnostic groups \citep{Poweretal2012}. In this study, head motion is measured using mean framewise displacement, which represents the average movement of the head between consecutive image volumes. Greater head motion is observed among autistic participants in the dataset (Table~\ref{tab:participants}), indicating that reported functional-connectivity differences may also be affected by motion. This issue is relevant to reported reductions in connectivity within the somatomotor system, which includes regions involved in movement and bodily sensation \citep{Holigaetal2019,Ilioskaetal2023}. Only the linear effect of motion is removed through motion regression. It therefore remains unclear whether this connectivity difference is preserved when autistic and non-autistic participants are explicitly matched for head motion.

Model design provides a second source of uncertainty. Increasingly complex nonlinear classifiers are proposed for autism classification, but their apparent gains are difficult to isolate when feature representations, cohorts, and validation procedures also differ. More complex nonlinear models are increasingly proposed for ASD classification \citep{Liuetal2024MADE,Tangetal2025GNNLSTM,Tangetal2026CASGNN}. Strong performance with tangent-space connectivity and regularised linear classifiers is also reported in an earlier benchmark \citep{Dadietal2019}. Regularised linear classifiers on tangent-space connectivity remain strong baselines despite the increasing use of nonlinear models. The independent effect of classifier choice cannot be determined reliably from comparisons across existing studies because other parts of the analysis pipeline also differ. Its effect is therefore evaluated using the same participants, features, and cross-validation folds. Performance alone also does not establish what the classifier score represents. Associations with symptom severity may reflect diagnostic-group separation, while apparent gains from windowed connectivity may arise from the stronger covariance regularisation required for short windows. Classifier scores have been reported to correlate with symptom severity and have been interpreted as evidence of a graded phenotype \citep{Nielsenetal2013}. When this correlation is calculated across both diagnostic groups, much of the association may be produced by group separation because severity measures also differ between the groups. It is therefore not established that classifier scores vary with symptom severity within autistic participants. Second, the contribution of temporal information beyond static connectivity is examined. Windowed connectivity is proposed to represent changes in connectivity during a scan \citep{Hutchisonetal2013,Allenetal2014,Pretietal2017}. Covariance estimates from short windows are known to be unstable and require additional regularisation \citep{LeonardiVanDeVille2015}. Additional regularisation may therefore contribute to the apparent advantage of windowed connectivity. For a fair comparison, the same regularisation is therefore applied to the whole-scan representation.

Two further requirements concern the reliability of the fitted model itself: its probabilities require calibration and its explanations require reproducibility. Hence, at first, the reliability of the probabilities returned by the ensemble is considered. Correct ranking of participants is possible even when predicted probabilities are systematically too high or too low. This property is not captured by the discrimination measures commonly reported in ABIDE-I benchmarks \citep{Guoetal2017,VanCalsteretal2019}. Second, the evidence used by the model is examined for interpretation at the level of brain systems. In a linear model, each connection weight is estimated together with all other connections and is not interpreted directly as a group difference \citep{Haufeetal2014}. Direct interpretation is made particularly unreliable when thousands of correlated connections are included. Both questions are therefore examined for the same ensemble to determine how reliably its scores and learned patterns can be interpreted.


% Part 3: existing solutions
Most connectome-based autism studies emphasise improvements in discrimination through alternative connectivity representations, deep models, graph methods, or site harmonisation. Comparatively less attention has been given to whether measured performance changes across participant subgroups, whether predicted probabilities are calibrated, and whether the patterns used to interpret a classifier reproduce in independent data. Different representations of functional connectivity, including the tangent-space projection of covariance matrices \citep{Varoquauxetal2010}, are compared in a controlled benchmark \citep{Dadietal2019}. Accuracies close to 0.70 are reported for deep autoencoder models on ABIDE-I \citep{Heinsfeldetal2018,Eslamietal2019,AlmuqhimSaeed2021}. Graph-based and convolutional models are also used to represent relationships between participants \citep{Parisotetal2018} and connections within brain networks \citep{Kawaharaetal2017}. These modelling approaches are further extended through multi-atlas deep ensembles \citep{Liuetal2024MADE}, dynamic graph models \citep{Tangetal2025GNNLSTM}, and multimodal graph models \citep{Tangetal2026CASGNN}. Differences between acquisition sites are addressed through harmonisation methods \citep{Johnsonetal2007,Fortinetal2018} and validation on held-out sites \citep{Abrahametal2017}. Performance in most of these studies is reported for the whole cohort, so variation in classification performance across participant groups is not clearly established. It also remains unclear whether reported differences in functional connectivity are preserved when head motion is directly controlled. In addition, the effect of classifier family is rarely examined under the same cross-validation setting, apart from the controlled benchmark reported \citep{Dadietal2019}. Together, these limitations motivate a controlled analytical framework in which the same connectome pipeline is used to separate the effects of participant composition, measurement confounding, modelling choices, probability calibration, and explanation stability.

% Part 4: what is proposed
To isolate these sources of variation, a fixed connectome-classification pipeline is used as an experimental framework rather than as the primary object of comparison. Three functional-connectivity representations are examined for each participant: correlation, static tangent-space connectivity, and windowed tangent-space connectivity, with the windowed representation derived from shorter segments of each scan. All three representations are analysed within the same regularised linear framework and evaluated using pooled ten-fold and leave-one-site-out validation, allowing within-cohort discrimination to be considered alongside transfer across acquisition sites. Discrimination is summarised using AUROC from held-out predictions. The framework is then used to examine how participant characteristics and head motion influence measured performance and connectivity differences, how modelling and regularisation choices affect the reported results, and whether classifier probabilities and explanations can be interpreted reliably. Functional-connectivity group contrasts are analysed independently of the classifier, while additional analyses are classified as primary descriptive, secondary, or exploratory according to their purpose and timing. From this framework, five main contributions are developed concerning probability calibration, reproducible model interpretation, motion-matched connectivity differences, the meaning of classifier scores across cognitive ability and symptom severity, and the effects of model complexity and covariance regularisation.

The main contributions are as follows:

\begin{itemize}

\item \textbf{Training-only calibration of connectome-classifier probabilities.}
Cross-fitted recalibration within each training partition removes the average probability offset without changing discrimination, improving calibration-in-the-large from $-0.162$ to $-0.010$. Residual calibration differences across withheld sites are strongly associated with the observed site-level autistic proportion, showing that average calibration is restored while site-specific probability shifts potentially remain. Thresholds selected from training data meet prespecified sensitivity and specificity targets in participant-pooled held-out predictions under both validation protocols, although these results are interpreted within the observed autistic proportion of the study cohort rather than as evidence for screening use.

\item \textbf{Reproducible interpretation of connectome classifiers at the level of brain systems.}
Each linear ensemble member is decomposed exactly into contributions from 36 functional-system pairs plus one unassigned connection group, and the resulting patterns are tested on non-overlapping participant halves within ABIDE-I. System-level attribution is reproducible across random non-overlapping halves (Spearman correlations 0.60--0.70) and remains stable when the halves share no acquisition site (0.53--0.63), whereas individual-connection patterns are substantially less stable. These results support interpretation at the functional-system level rather than at the level of isolated connections.

\item \textbf{Motion-matched evidence for somatomotor connectivity differences and system-level classifier attribution.}
The somatomotor reduction strengthens from $-0.215$ to $-0.249$ after within-site matching on recorded head motion, and widespread reductions remain under residual permutation. In the classifier, default-mode pairs account for the largest overall share of group separation, while the somatomotor system contributes most per connection. Agreement between regional classifier attribution and regional group differences is also observed within the somatomotor system. These findings are treated as exploratory associations rather than biomarkers.

\item \textbf{Characterisation of what a connectivity-based autism score reflects across cognitive ability and symptom severity.}
Measured discrimination varies with the cognitive composition of the autistic sample. In unmatched IQ bands, AUROC decreases from 0.873 to 0.719, while diagnosis-by-IQ analyses show stronger negative score--IQ associations in autistic participants for full-scale and verbal IQ. In contrast, symptom associations observed across diagnostic groups do not survive correction within autistic participants. The classifier score is therefore supported as a measure of diagnostic-group separation whose performance depends on cognitive composition, rather than as a measure of autism severity.

\item \textbf{Controlled evidence on model complexity and covariance regularisation in connectome classification.}
Under identical features and cross-validation folds, regularised linear configurations achieve AUROCs of 0.735--0.757, compared with 0.611--0.624 for fixed-configuration tree ensembles. The apparent advantage of windowed connectivity is also removed when the same covariance shrinkage is applied to the whole-scan representation, leaving an AUROC difference of $-0.002$. These results support regularised tangent-space connectivity as a strong baseline for evaluating more complex classifiers and temporal representations.


\end{itemize}
The remainder of the paper is organised as follows: related work is reviewed in
Section~\ref{sec:related}, the methods are described in
Section~\ref{sec:methods}, and the results are reported in Section~\ref{sec:results}.
The findings are then discussed in Section~\ref{sec:discussion}, and conclusions are drawn in
Section~\ref{sec:conclusion}.
